% !PRINTED 692
\weberpage{692}
%<<CONT
it follows that, if $\vartheta$ is taken within the limits
\[
0 \leqq \vartheta < 2\pi
\]
then, for a fixed $z$, there belongs to every value of $\vartheta$ a corresponding pair of values $u_1, u_2$. Such an angle $\vartheta$ belongs to every imaginary pair.

To every point of the region $S$ there then belongs one and only one system of values of the variables $\xi_1, \xi_2, \ldots, \xi_\nu, \vartheta \ldots$ within the bounds

\weq{7}{\begin{aligned}
&0 \leqq \xi_1 < 1, \ldots, 0 \leqq \xi_{\nu-1} < 1 \\
&\xi_\nu < 0, \qquad 0 \leqq \vartheta < 2\pi, \ldots
\end{aligned}}
% [source notation] Both first-row less-than-or-equal signs place the equality bars above the angle; other occurrences on this page place them below. Both variants express the same relation and are represented by \leqq. No damaged-type or reversed-inequality claim is made.

If conversely any system of values of the variables $\xi, \vartheta$ within the limits (7) is given, one obtains from it the corresponding parts $u_1, u_2$ of an imaginary $y$ uniquely, but of the real $y$ only the absolute values, and so each of the real $y$ may still be given either of the two signs. The number of these possible determinations is $2^{2\nu-n}$. If one of these determinations has been singled out, the variables $x$ are uniquely determined, and they are finite for every finite system of values of the $\xi, \vartheta$. For a $\xi_\nu$ decreasing towards $-\infty$ the values of the $y$ will in part converge to zero.

The region $S$ accordingly splits into $2^{2\nu-n}$ partial regions $S_1, S_2, \ldots$, each of which is characterised by a definite sign-combination of the real $y_1, y_2, \ldots$, and accordingly the integral $U$ splits into just as many partial integrals $U_1, U_2, \ldots$

\[
U = U_1+U_2+\cdots
\]

In carrying out the integration in one of these parts, all of which prove to be of equal value, we begin with an imaginary pair, if such a one is present.

If we regard $u_1, u_2$ as rectangular coordinates in a plane, then by (5) we may conceive $e^{1/2\,z}, \vartheta$ as polar coordinates, and obtain

\[
\int\int du_1\,du_2 = \frac12\int\int e^z\,dz\,d\vartheta,
\]

in which the integration with respect to $\vartheta$, which extends from $0$ to $2\pi$, can be carried out. One thus finds

\[
\int\int du_1\,du_2 = \pi\int e^z\,dz.
\]
